\section{Заключение}

В ходе лабораторной работы были изучены выпуклые множества, выпуклые линейные комбинации, выпуклые оболочки и свойство сохранения выпуклости при пересечении. Теоретические положения были проверены вычислительными экспериментами.

В первой задаче для случайного набора точек была построена выпуклая
оболочка. Полученный многоугольник содержит все исходные точки и
является выпуклым по построению. 

Во второй задаче были заданы два
выпуклых круга, а их принадлежность точкам плоской сетки
визуализирована. Форма общей части и теорема о пересечении
подтверждают, что пересечение выпуклых множеств выпукло.

В третьей
задаче реализована проверка выпуклости многоугольника через сравнение
знаков ориентированных произведений: тесты дали ожидаемые результаты
для выпуклой, невыпуклой и звезднообразной фигур соответственно.

Ограничение реализованного критерия состоит в том, что он рассчитан на упорядоченный список вершин простого многоугольника. Для произвольного неупорядоченного набора точек сначала следует построить выпуклую оболочку, а для самопересекающихся контуров требуется отдельная проверка корректности границы.

Использование выпуклости позволяет надежно описывать области допустимых решений и применять эффективные методы оптимизации: если задача является выпуклой, поиск
локального минимума приводит к глобально оптимальному решению. При
этом несколько выпуклых ограничений можно объединять, поскольку их
пересечение также остается выпуклым.

Таким образом, поставленная цель лабораторной работы была достигнута:
основные свойства выпуклых множеств были рассмотрены теоретически и
подтверждены с помощью программных экспериментов. Это показывает, что выпуклость является
не только теоретическим свойством геометрических объектов, но и
удобным инструментом для анализа ограничений и решения задач
оптимизации.
